\section{A graph-driven recursive BPH Gaussian compiler}
\label{sec:automatic-compiler}

The input is a connected, amputated, one-particle-irreducible multigraph
$G$ of massive scalar $\phi^3$ theory in six Euclidean dimensions.
Internal edges retain distinct indices, including parallel edges; external
half-edges retain their physical labels.  Let $L$, $I$ and $E$ denote its
loop, internal-edge and external-leg counts.  Trivalence and Euler's
identity give
\[
 I=3L+E-3,\qquad \Omega_G=6L-2I=6-2E.
\]
The present implementation accepts $1\le L\le4$ and $1\le E\le3$.
Its generated expressions use no topology-specific integration kernel.
This section establishes the regulated construction and its interior
parameter formula, and states separately the analytic and numerical
evidence needed to identify and integrate its unregulated limit.

\subsection{The primary construction is the Bogoliubov recursion}

For an embedded edge subgraph $\gamma$, all omitted incident half-edges
are separate boundary ports.  This convention matters when parallel
edges have the same endpoints.  The compiler enumerates proper,
connected, 1PI edge subsets with
\[
 \omega_\gamma=6L_\gamma-2I_\gamma\ge0.
\]
There is no replacement of an edge subset by an induced vertex subset.
If $P_\gamma$ lists all independent boundary momenta, the Taylor operator
is the full total-degree polynomial
\[
 t_\gamma^{\omega_\gamma} f(P_\gamma)
 =\sum_{r=0}^{\omega_\gamma}\frac1{r!}
   \left.\frac{d^r}{dz^r}f(zP_\gamma)\right|_{z=0}.
\]
The boundary variables can include other loop momenta of $G$; they are
not restricted to the physical external momenta.

For a nonempty spinney $S$ of mutually vertex-disjoint proper divergent
subgraphs, put
\begin{align*}
 \overline R_G I_G
 &=I_G+\sum_{S\ne\varnothing}
       \left(\prod_{\gamma\in S} C_\gamma\right)I_{G/S},\\
 C_\gamma&=-t_\gamma^{\omega_\gamma}\overline R_\gamma I_\gamma,
 \qquad
 R_G I_G=(1-t_G^{\Omega_G})\overline R_G I_G.
\end{align*}
The local polynomial carried by a two-point counterterm is retained in
full, including its mass and kinetic terms.  Contracting a subgraph in
this notation does not discard that polynomial.

The actual program expands this recursion by choosing disjoint maximal
subgraphs and then recursively preparing their proper subgraphs.  It
records a word $W$ with its subgraphs ordered from inner to outer and
coefficient $(-1)^{|W|}$.  The maximal members and the successive maximal
members inside them are unique for any resulting word.  There is
therefore no additional multiplicity attached to a word.  For the
finite input graph the resulting support can be compared with a
separately enumerated set of compatible forests.  This is an auxiliary
finite combinatorial check; the implementation does not invoke an
unverified general forest or $Z$ identity to construct its answer.

\subsection{A compatible tree gives exact integer momentum coordinates}

For a fixed recursive word, its selected subgraphs are nested or
vertex-disjoint.  A spanning tree can be constructed inside each minimal
selected subgraph, extended inside its parent, and finally extended to
$G$.  Equivalently, the implemented Kruskal ordering first processes
edges belonging to the smallest selected subgraph, breaking ties by the
original edge index.  For every selected $\gamma$,
$T\cap E_\gamma$ is then a spanning tree of its vertices.

Let the edges outside $T$ be the chords $h_1,\ldots,h_L$.  Assigning a
unit momentum to each chord and solving conservation along tree leaves
gives
\[
 k_e=\sum_{j=1}^L r_{ej}\ell_j
       +\sum_{a=1}^{E-1}r_{e,L+a}p_a,
 \qquad r_{ej}\in\mathbb Z,
 \qquad p_E=-\sum_{a<E}p_a.
\]
The reduced oriented incidence matrix of a tree has determinant $\pm1$.
Consequently this solve is integral, and different fundamental cycle
coordinates are related by an integer unimodular transformation,
possibly followed by a shift linear in the physical momenta.  Its loop
integration Jacobian is one.  An internal chord of $\gamma$ has its tree
path entirely in $\gamma$, so its fundamental cycle vanishes on every
edge outside $\gamma$.  The compiler checks both the number of these
chords, $L_\gamma$, and this absence of leakage.

At a positive common Schwinger cutoff, loop integration and Taylor
differentiation may be interchanged.  Translation invariance of those
Gaussian integrations also makes the integrated local Taylor polynomial
independent of the compatible routing used for its word.  This does
not assert that an external-momentum-dependent loop shift commutes
pointwise with an unintegrated Taylor operator.

\subsection{Scaled rows implement the boundary Taylor operators}

For each selected $\gamma$ introduce a formal scalar $t_\gamma$.  For
each edge $e$ and each coordinate $j$ of the routing, define
\[
 \chi_\gamma(e,j)=
 \begin{cases}
 1,&e\in E_\gamma\ \text{and }j\text{ is not an internal chord of }\gamma,\\
 0,&\text{otherwise}.
 \end{cases}
\]
Every physical coordinate is external to $\gamma$ in this definition.
The row used to build its Gaussian form is
\[
 c_{ej}(t)=r_{ej}\prod_{\gamma\in W}
                 t_\gamma^{\chi_\gamma(e,j)}.
\]
Thus a boundary loop momentum is scaled even though it is not a physical
external momentum.  In a nested pair an inner physical coefficient can
contain $t_{\rm inner}t_{\rm outer}$, an outer internal loop coefficient
inside the inner subgraph contains only $t_{\rm inner}$, and an inner
internal loop coefficient is unchanged.  These are precisely the
inside-out boundary scalings required by the recursion.

For $\omega_\gamma=0$ the program sets $t_\gamma=0$ directly.  For the
other selected subgraphs it works in the finite ring
\[
 \mathbb R[t_\gamma:\gamma\in W,\ \omega_\gamma>0]
     /(t_\gamma^{\omega_\gamma+1}:\gamma\in W),
\]
and finally sums all coefficients in the rectangular index set
$0\le k_\gamma\le\omega_\gamma$.  A coefficient in this ring already
means $\partial^k f(0)/\prod k_\gamma!$.  Multiplying it by another
factorial would be an error.  The rectangular indexing concerns the
different Taylor operators; each individual operator is still a
total-degree Taylor polynomial in all of its boundary momenta.

\subsection{Gaussian integration in the finite Taylor algebra}

Give every original edge its own positive Schwinger parameter $\alpha_e$.
No parameter is removed when a recursive word inserts a counterterm.
For that word set
\[
 Q_W(t)=\sum_{e=1}^I\alpha_e c_e(t)c_e(t)^T
       =\begin{pmatrix}A&B\\B^T&C\end{pmatrix},\qquad
 U_W=\det A,
\]
where $A$ is the $L\times L$ loop block.  With the physical Gram matrix
$g_{ab}=p_a\cdot p_b$, define
\[
 H_W=\sum_{a,b=1}^{E-1}g_{ab}
                  (C-B^TA^{-1}B)_{ab}.
\]
The constant loop block is positive definite at an interior parameter
point: the internal subgraph cycles and the quotient cycles still form
a complete basis after its selected boundary momenta have been set to
zero.  Gaussian elimination is therefore legitimate in the Taylor
ring.  Inverting a pivot uses a finite geometric expansion in its
nilpotent remainder; a logarithm uses the corresponding finite Taylor
series.  These algebraic operations implement derivatives of the
Gaussian integral, rather than numerical finite differences.

For each of the six spatial coordinates the same scalar loop matrix
occurs.  With the convention $d^6\ell/\pi^3$ per loop,
\[
 \int\prod_{j=1}^L\frac{d^6\ell_j}{\pi^3}
       e^{-\sum_e\alpha_e k_e(t)^2}
       =U_W(t)^{-3}e^{-H_W(t)}.
\]
There is no diagram symmetry weight, coupling, $(4\pi)$ factor, or
dimensional-regularization prefactor in this equation.

\subsection{The overall subtraction makes the scale integral explicit}

Write $\alpha_e=\rho x_e$, $\sum_e x_e=1$.  The Jacobian is
$\rho^{I-1}d\rho\,d^{I-1}x$, $U(\rho x)=\rho^LU(x)$, and
$H(\rho x)=\rho H(x)$.  Hence the remaining scale power is
\[
 \rho^{I-1-3L}=\rho^{E-4}.
\]
The simplex measure here is the usual coordinate Lebesgue measure,
whose volume is $1/(I-1)!$.  No extra simplex factorial is to be inserted.

Since $H$ is quadratic in all physical momenta, the overall zero-degree
subtraction for $E=3$ removes $1$ from $e^{-\rho H}$; the full
degree-two subtraction for $E=2$ removes $1-\rho H$.  For $M=m^2>0$,
\begin{align*}
 \Phi_3(H,M)
 &=\int_0^\infty\frac{e^{-M\rho}}\rho
                    (e^{-H\rho}-1)d\rho
 =-\log(1+H/M),\\
 \Phi_2(H,M)
 &=\int_0^\infty\frac{e^{-M\rho}}{\rho^2}
                    (e^{-H\rho}-1+\rho H)d\rho\\
 &=(M+H)\log(1+H/M)-H.
\end{align*}
For $E=1$ momentum conservation leaves no independent external momentum;
the full zero-momentum Taylor condition therefore gives $R_GI_G=0$.
This includes the chosen finite condition on the linear counterterm.

For a word define $\mathcal T_W$ to set its zero-order variables to zero
and sum the retained Taylor coefficients of its other variables.
The projective density emitted by the compiler is
\[
 \mathcal D_G(x;p,M)=
   \sum_{W\text{ from the recursion}}(-1)^{|W|}
    \mathcal T_W\bigl[U_W(t)^{-3}\Phi_E(H_W(t),M)\bigr].
\]
All words are evaluated at the same original $x$ and combined before
numerical integration.  In particular a separately divergent word is
not assigned an unregulated numerical integral.

\subsection{A two-loop overlapping graph extracted from the actual program}

The following matrices were extracted by
\path{tools/extract_automatic_overlap.py} from actual
\texttt{compile\_bph} output, retained in
\path{docs/automatic-compiler-evidence/generated-overlap.json}.  They are an illustration of the
algorithm, not a formula table used by its engine.  The edge order is
\[
 (0,2),(0,3),(1,2),(1,3),(2,3),
\]
with physical legs at vertices $0,1$.  Write its parameters as
$(a,b,c,d,e)$ and $p^2=s$.  Its two proper divergent triangles have edge
sets $\gamma_a=\{0,1,4\}$ and $\gamma_b=\{2,3,4\}$.
They overlap, so the recursive words are only
$\varnothing$, $\{\gamma_a\}$ and $\{\gamma_b\}$.

For the empty word, the program obtains the routing rows in the basis
$(\ell_1,\ell_2,p)$
\[
 \begin{pmatrix}1&1&1\\-1&-1&0\\-1&0&-1\\1&0&0\\0&1&0\end{pmatrix}.
\]
The full loop-plus-physical matrices after each selected zero-order
Taylor substitution are
\[
 Q_\varnothing=
 \begin{pmatrix}
 a+b+c+d&a+b&a+c\\
 a+b&a+b+e&a\\
 a+c&a&a+c
 \end{pmatrix},
\]
\[
 Q_{\gamma_a}=
 \begin{pmatrix}c+d&0&c\\0&a+b+e&0\\c&0&c\end{pmatrix},
 \qquad
 Q_{\gamma_b}=
 \begin{pmatrix}a+b&0&-a\\0&c+d+e&0\\-a&0&a\end{pmatrix}.
\]
The compatible tree and chord basis of the last word differ from those
of the first word; each transformation has unit Jacobian.  Put
\begin{align*}
 U&=(a+b)(c+d)+e(a+b+c+d),\\
 V&=ab(c+d)+cd(a+b)+e(a+c)(b+d).
\end{align*}
Direct determinants and Schur complements of these emitted matrices give
\begin{align*}
 U_\varnothing&=U,&H_\varnothing&=sV/U,\\
 U_{\gamma_a}&=(c+d)(a+b+e),&H_{\gamma_a}&=s\,cd/(c+d),\\
 U_{\gamma_b}&=(a+b)(c+d+e),&H_{\gamma_b}&=s\,ab/(a+b).
\end{align*}
The resulting density is the first $U^{-3}\Phi_2$ expression minus
the other two.  In particular no simultaneous
$t_{\gamma_a}t_{\gamma_b}$ term is introduced for the overlapping pair.

\subsection{A common regulator and the exact scope of the limit}

For regulated algebra use the same propagator on every original edge,
\[
 D_\eta(k)=\int_\eta^\infty e^{-\alpha(k^2+M)}d\alpha,
 \qquad \eta>0.
\]
Every recursive word retains the same original region
$\alpha_e\ge\eta$.  All Gaussian integrals and finitely many Taylor
derivatives are then absolutely integrable.  After projectivization the
common lower scale limit is
\[
 T(x,\eta)=\eta/\min_e x_e.
\]
Replacing the lower limits in the two definitions of $\Phi_E$ by this
same $T$ gives the regulated Gaussian expression exactly.  For fixed
interior $x$ the subtracted exponential is $O(\rho)$ for $E=3$ and
$O(\rho^2)$ for $E=2$, also after any of the finitely many Taylor
derivatives.  The mass controls infinity.  It follows that its
$\eta\downarrow0$ limit is the displayed projective density.

That interior conclusion alone does not justify taking the limit
through the entire simplex or through the original $6L$-dimensional
momentum integral.  To make both identifications one needs suitable
integrable majorants for the complete recursive sum, uniform in the
common regulator.  A future inductive proof must also retain control in
the internal subgraph-external momenta, not just at one fixed reference
scale.  The separate one- and two-loop family proofs provide this
bridge for their stated domains.  Graph-derived primitive certificates
provide additional bounds for their explicitly certified graphs.
Compiling all three- and four-loop
graphs, or testing a collection of UV rays, is not a proof of the
corresponding uniform bridge for every such graph.  The compiler's
formal expression, a pointwise numerical UV diagnostic, an integrated
estimate, and an analytic convergence certificate are distinct claims.

\subsection{Arithmetic cancellation is checked separately from integration}

The common-parameter sum can be much smaller than its individual words.
For the two-loop nested bubble with edges
\[
 (0,1),(0,3),(1,2),(2,3),(2,3)
\]
and $s=M=1$, an actual randomized-integration point was
\[
 \begin{aligned}
 \alpha={}&(0.6311634306896491,\ 0.016292580874392757,\ 0.35254392304634735,\\
 &\qquad6.228782654682385\!\times10^{-8},\
 3.101784279669877\!\times10^{-9}).
 \end{aligned}
\]
Its unguarded binary64 density was $+440112873472$, while independent
100-digit differentiation gave $-3680.2843093271767$.  A positive
importance weight of about $1.71\times10^{-8}$ turned the false density
into a spurious weighted value about $7527$, rather than the correct
$-6.29\times10^{-5}$.  This was an arithmetic error, not statistical
evidence for a large positive tail.  A signed local BPH density is allowed
and is not clipped to impose positivity.

The initial safeguard recomputed the complete sum whenever a
proper-subgraph parameter sum or its complement was below $10^{-4}$.
That face check is retained, but simultaneous quadratic subtractions
can fail outside its threshold. The current whole-coefficient relative
trigger and optional checked native executor are described in
Section~\ref{sec:precision-repair}; both retain the independent standard
library Decimal path. If the smallest positive face scale, also
including $1$, is $r$, the Decimal starting precision is
\[
 d=\min\{470,\max\{80,32+12\lceil\max(0,-\log_{10}r)\rceil\}\}.
\]
It then recomputes with $\min(500,d+30)$ digits, requiring
\[
 |v_{d'}-v_d|\le10^{-25}\max\{|v_{d'}|,|v_d|\},
 \qquad d'=\min(500,d+30),
\]
with an explicit acceptance when both values are zero, and increases
precision again if necessary. Failure to stabilize by
500 digits raises an error. The displayed failing point now agrees
with the independent reference.  The batch path tries
\texttt{numpy.longdouble} first, but records its actual mantissa width;
on the tested platform it has only 53 significant bits.  No extra
precision is inferred from its type name.

Increasing the arithmetic precision was not by itself sufficient.  A
second test exposed a small-$z$ series truncated at $n=29$ for
$\Phi_2$, with $z=H/M$, that had been reused even in high-precision
arithmetic.  For parameters
\[
 (0.1,0.4,0.5,4\times10^{-21},6\times10^{-21}),\quad s=M=1,
\]
independent 160-digit differentiation gave
\[
 116357167448485560.1524657840956468920609\ldots,
\]
whereas both 110- and 200-digit calculations with that fixed series
stabilized at
\[
 116357167448126952.4750081595388072343536\ldots.
\]
Their common bias escaped an arithmetic precision comparison.  The
omitted analytic tail, before multiplication by $U^{-3}$, is bounded by
$M|z|^{30}/[30\cdot29(1-|z|)]$ when $|z|<1$.  It can therefore be
amplified near a UV face.  The Decimal implementation now uses
context-dependent convergent summation for both small-$z$ functions:
\[
 \log(1+z)=\sum_{n\ge1}\frac{(-1)^{n+1}z^n}{n},\qquad
 \Phi_2=M\sum_{n\ge2}\frac{(-1)^nz^n}{n(n-1)}.
\]
It continues until adding the next term no longer changes the sum in
the active context, then performs the higher-precision recheck.  Both
failure cases are regression tests, rather than being hidden by a
wider numerical tolerance.

Decimal recomputation evaluates the accepted binary64 input parameters
and Gram entries; it does not invent extra input accuracy. The face
threshold, whole-coefficient condition trigger and precision comparisons
are practical arithmetic diagnostics,
not outward-rounded error intervals or a theorem that all cancellation
patterns have been detected.

\subsection{Numerical integration and current verification}

The optional generic integrator maps independently scrambled Sobol
points to Gamma variables and normalizes them to a symmetric Dirichlet
distribution on the original $I$-edge simplex.  For parameter $\beta$,
\[
 p_\beta(x)=\frac{\Gamma(I\beta)}{\Gamma(\beta)^I}
                   \prod_e x_e^{\beta-1},\qquad
 J=\mathop{\mathrm E}_{p_\beta}
                   [\mathcal D_G(x)/p_\beta(x)].
\]
The full word sum is formed before multiplication by this importance
weight and before averaging.  Standard errors use the variation of the
independently scrambled replicate means, not an independence assumption
for the points within one scramble.  The first half of each nested
Sobol sequence supplies a second power diagnostic.  These quantities,
maximum weighted samples, native precision checks and actual Decimal
fallback counts are recorded
separately; none is a rigorous total quadrature enclosure.

All 698 generated one- through four-loop graphs compile.  Their recursive
supports agree with the independent finite forest-support check.
Independent exact rational Gaussian determinants agree with
matrix-tree and spanning-two-forest polynomials for the tested original
words, including selected three- and four-loop primitive graphs.
Closed one-loop finite densities, physical-label and edge/vertex
permutations, scalar/batch agreement, mass dimension, zero subtraction
points, the two historical cancellation counterexamples and the new
simultaneous-face cases in Section~\ref{sec:precision-repair} are tested. Old
topology-specific entry points are patched to raise during an additional
test.  These checks substantiate the implemented compiler and its
recorded numerical results; they do not close an arbitrary-loop BPH
convergence induction or certify all three- and four-loop numerical
integrals.
